\chapter{Pilot Design}

The pilot is the most important commercial moment in 2026.
Getting the pilot structure right matters more than getting the product perfect.

The goal of the pilot is not to show everything Connector can do.
The goal is to prove one thing clearly:

\begin{conceptbox}{Pilot goal}
\centering
\large\bfseries After this pilot, the customer can see what their AI did and why.
\end{conceptbox}

\section{Core design principle --- one workflow}

The right pilot is:

\begin{itemize}
  \item one workflow,
  \item one team,
  \item one operational problem,
  \item one success definition.
\end{itemize}

Not multi-team.
Not a broad rollout.
Not ``let us instrument the whole company.''

The narrower the pilot, the faster it succeeds.
The faster it succeeds, the easier expansion becomes.

\section{What we show in the pilot}

For the chosen workflow, we show:

\begin{enumerate}
  \item \textbf{Full trace.} Every step the AI took, in order, with data context.
  \item \textbf{Explanation of the decision path.} Why the AI reached the output it reached.
  \item \textbf{Audit-ready output.} Something a non-technical stakeholder can review and understand.
  \item \textbf{Clearer operational visibility} than the team has today without Connector.
\end{enumerate}

The customer should come away with a simple reaction:

\begin{insightbox}{The reaction we want}
\centering
\large ``Now we can see what happened.''
\end{insightbox}

That reaction is more valuable than any technical feature demonstration.

\section{Pilot timeline}

\begin{center}
\renewcommand{\arraystretch}{1.5}
\begin{tabularx}{0.88\textwidth}{@{}lX@{}}
\toprule
\textbf{Context} & \textbf{Timeline} \\
\midrule
Fast-moving startup or internal team & 2--3 weeks \\
Mid-size company with some review process & 3--4 weeks \\
Larger org with security or procurement review & 4--6 weeks \\
\bottomrule
\end{tabularx}
\end{center}

The principle is the same regardless of timeline: short, narrow, concrete.

\section{What the pilot is not}

\begin{itemize}
  \item It is not a proof of concept for the whole platform.
  \item It is not a free engagement.
  \item It is not indefinite.
  \item It is not a replacement for a commercial agreement.
\end{itemize}

The pilot has a defined scope, a defined timeline, and a defined success criteria.

\section{What success looks like}

A pilot is complete and successful when the customer can say:

\begin{enumerate}
  \item We can reconstruct what the AI did inside this workflow.
  \item We can explain why it reached that output.
  \item We can show this evidence to internal stakeholders.
  \item We want this on at least one more workflow.
\end{enumerate}

The fourth point is the expansion bridge.
That is where a pilot becomes a commercial relationship.

\section{Pricing model for pilots}

The pilot should be paid.
A free pilot signals low confidence and attracts low-quality evaluation.

A reasonable pilot price is between \$2,000 and \$8,000 depending on the segment:

\begin{center}
\renewcommand{\arraystretch}{1.5}
\begin{tabularx}{0.88\textwidth}{@{}lX@{}}
\toprule
\textbf{Segment} & \textbf{Pilot price range} \\
\midrule
AI support startup & \$2,000--\$4,000 \\
Fintech ops team & \$4,000--\$8,000 \\
Internal automation at a mid-market company & \$3,000--\$6,000 \\
\bottomrule
\end{tabularx}
\end{center}

A paid pilot also creates a clear transition path to a monthly or annual contract if it succeeds.
